\section{Conclusion}

James Wilkinson survived because he was useful before he was defensible. Jefferson's relationship with him was built on that usefulness, especially in the western theater and during the Burr crisis. When accusations mounted, Jefferson's support and later Madison's approval did not erase suspicion. They did something narrower and more political: they kept him within reach of formal vindication. In that respect Wilkinson belongs to a larger early republic story in which reputation, command, and executive management were closely entangled \citep{freeman2002,kohn1975,stagg1983}.

Frederick matters because viability needed a place. Wilkinson's older Maryland footing and the barracks' connection to Jeffersonian western logistics still matter as background, but the heart of the case now lies elsewhere: the 1800 War Department receipt, the 1804 Frederick letter to Adams, Jefferson's 1808 public insistence on inquiry, the 1811 venue and witness letters between Eustis and Madison, the interrogatories sent from Frederick Town, Wilkinson's own plea for an unbiased hearing, Adams's testimonial reply, and the 1812 approval of the Frederick-town acquittal. Read together, those records show Frederick as one of the towns in which national scandal became manageable federal procedure \citep{lurie2001}.

That is the paper's main contribution. It reframes a notorious national figure through the local machinery that made his survival possible. For a Frederick historical society, the most important takeaway is not that Wilkinson once appeared there. It is that Frederick helped turn a national scandal into a sequence of proceedings that could end in acquittal and restored standing. Frederick was not an incidental stage set. It was one of the working mechanisms.

The note should still be read with evidentiary discipline. The core relationship claim rests most securely on direct correspondence and official records. The Frederick claim is now stronger than a simple local gloss because most of its documentary chain is primary, but the thicker local texture still depends in part on a small number of later Frederick-centered sources. That does not weaken the synthesis, but it does define its proper scale: this paper offers a strong, source-accountable interpretation, not the last archival word on Wilkinson.

That scale is also what makes the paper useful for public history. A local historical society does not need Wilkinson to become admirable in order for Frederick to matter. It needs a documented account of how local place entered a national political story. On that narrower but important question, the evidence is strong. Frederick appears not as a sentimental hometown cameo but as one of the places through which executive protection was turned into durable official form.

There is also a public-history reason to keep difficult actors in view. Wilkinson is not valuable because he is morally exemplary. He is valuable because compromised figures often show most clearly how institutions, places, and reputations actually work. Following him does not ask Frederick to celebrate him. It asks the town to recognize that local history can illuminate national power most sharply when it traces how messy lives were processed, defended, and judged.

If this note were expanded for a broader scholarly venue, the next archival step would be clear: a deeper integration of the printed court-martial report, the NARA Record Group 153 case file, the Jefferson Papers guides and correspondence lists, Frederick newspaper abstracts and surviving issues, Frederick legal manuscripts, and the Taney papers to test how far the Frederick-centered procedural story can be extended beyond the compact documentary chain used here \citep{morgancourtmartial1812,nara153,locjeffersonpapers,locjeffersonwilkinsonlist,locjeffersonwilkinson1811list,wright1992,uvataneypapers}.
